\chapter{Theoretical Background}
\label{chap:background}

This chapter introduces the theoretical concepts underlying the
implementation and empirical evaluation, including the online-learning
setting, regret notions, equilibrium concepts, and their main connections.


\section{Online Learning Setting}
\label{sec:online-learning-setting}

We consider a sequential decision problem with the finite action set
\([K]=\{1,\ldots,K\}\). At each round \(t=1,\ldots,T\), the learner chooses
an action without knowing the reward vector of that round in advance. Rather
than selecting an action deterministically, the learner may randomize by
choosing a probability distribution \(p_t\in\Delta_K\), where
\[
    \Delta_K
    =
    \left\{
        p\in\mathbb{R}_{\geq0}^{K} :
        \sum_{i=1}^{K}p(i)=1
    \right\},
\]
and then samples an action \(I_t\sim p_t\).

Randomization prevents the learner from always following a predictable
deterministic pattern and, under bandit feedback, also provides a mechanism
for exploring different actions.

Associated with round \(t\) is a reward vector
\(r_t=(r_t(1),\ldots,r_t(K))\in[0,1]^K\). After choosing \(I_t\), the learner
receives the reward \(r_t(I_t)\). Its objective is to obtain high cumulative
reward \(\sum_{t=1}^{T}r_t(I_t)\).

The reward sequence may be generated independently of the learner's realized
actions or may adapt to the previous interaction history. When rewards depend
on past interaction, the environment is assumed to be
\emph{non-anticipating}: \(r_t\) may depend on previous actions, but not on
the learner's fresh random draw \(I_t\) in the same round. The concrete reward
environments used in the experiments are described in
Chapter~\ref{chap:methodology}.

\paragraph{Feedback models.}
Under \emph{full-information feedback}, the learner observes the complete
reward vector \(r_t\) after each round.

Under \emph{bandit feedback}, the learner observes only the reward
\(r_t(I_t)\) of the sampled action. The rewards of the remaining actions are
hidden from the learner. EXP3 is a classical algorithm designed for this
limited feedback setting \citep{auer_nonstochastic_2002}.

The reward-generation process and the feedback model are separate aspects of
an online-learning problem. The experiments in this thesis consider both
full-information and bandit feedback.

Throughout the thesis, algorithms are described in terms of rewards. Results
stated in terms of losses can be translated using
\(\ell_t(i)=1-r_t(i)\).


\section{Regret Notions}
\label{sec:regret-notions}

Absolute cumulative reward alone is often insufficient to evaluate an online
learner, because the attainable reward depends strongly on the environment.
Regret instead compares the learner's decisions with alternative decisions
that could have been applied to the same reward sequence.

Throughout this thesis, regret computed from the actions actually sampled by
the learner is called \emph{action regret}, following the terminology of
\citet{greenwald_bounds_2006}. This is distinct from \emph{distribution
regret}, which is computed from the learner's mixed strategy rather than from
the sampled actions. For two actions \(i,j\in[K]\), define the cumulative
replacement gain of the sampled action sequence by
\[
    G_T(i,j)
    =
    \sum_{t=1}^{T}
    \mathbf{1}\{I_t=i\}
    \bigl(r_t(j)-r_t(i)\bigr).
\]
This quantity measures how much additional reward would have been obtained
if every realized occurrence of action \(i\) had instead been replaced by
action \(j\).

The matrix \(G_T\) is the common basis for all three regret notions used in
the experiments.

\begin{definition}[External regret]
The external regret after \(T\) rounds is
\[
    R_T^{\mathrm{ext}}
    =
    \max_{j\in[K]}
    \sum_{t=1}^{T}
    \bigl(r_t(j)-r_t(I_t)\bigr).
\]
Equivalently,
\[
    R_T^{\mathrm{ext}}
    =
    \max_{j\in[K]}
    \sum_{i=1}^{K}G_T(i,j).
\]
\end{definition}

External regret compares the realized trajectory with the best single fixed
action in hindsight.

\begin{definition}[Internal regret]
The internal regret after \(T\) rounds is
\[
    R_T^{\mathrm{int}}
    =
    \max_{i,j\in[K]}G_T(i,j).
\]
\end{definition}

Internal regret asks whether replacing one particular action \(i\) by another
action \(j\) would have systematically improved the learner's realized
performance \citep{hart_simple_2000}.

\begin{definition}[Swap regret]
A swap deviation is a mapping \(\phi:[K]\rightarrow[K]\) that assigns a
replacement action to each source action. The swap regret is
\[
    R_T^{\mathrm{swap}}
    =
    \max_{\phi:[K]\rightarrow[K]}
    \sum_{t=1}^{T}
    \bigl(
        r_t(\phi(I_t))-r_t(I_t)
    \bigr).
\]
Because the replacement for each source action can be chosen independently,
this is equivalently
\[
    R_T^{\mathrm{swap}}
    =
    \sum_{i=1}^{K}
    \max_{j\in[K]}G_T(i,j).
\]
\end{definition}

Swap regret therefore allows a separate best replacement for every action
selected by the learner \citep{goos_general_2003,blum_external_2007}. The
external-, internal-, and swap-regret quantities above are all action-regret
quantities because they are evaluated on the sampled action sequence.

\paragraph{Expected action regret.}
For a randomized learner, action regret is itself a random variable. For
\(q\in\{\mathrm{ext},\mathrm{int},\mathrm{swap}\}\), this thesis uses
\[
    \mathcal{R}_T^{q}
    =
    \mathbb{E}\!\left[R_T^{q}\right]
\]
for the corresponding \emph{expected action regret}. The expectation is over
the randomness of the interaction: the learner's randomization and, where
applicable, random environments or other randomized players. In the swap
case, this convention agrees with the ``expected swap regret'' studied by
\citet{huang_near-optimal_2023}, where the best swap deviation is selected
inside the expectation.

Huang and Pan distinguish this quantity from swap pseudo-regret, where the
maximization over deviations is taken outside the expectation
\citep{huang_near-optimal_2023}. The two notions are therefore not treated as
interchangeable in this thesis. The experiments estimate expected action regret
by first computing action regret separately for each replicate and then
averaging the resulting values.

\paragraph{Relations between the regret notions.}
The three notions differ in the deviation class used for comparison.
External regret considers one fixed replacement, internal regret one
source--target replacement, and swap regret an arbitrary replacement mapping.
Blum and Mansour note that external and internal regret are both
upper-bounded by swap regret, while swap regret is at most \(K\) times
internal regret; external and internal regret are not generally ordered
\citep{blum_external_2007}. Thus swap regret is the strongest of the three
benchmarks considered here, while external and internal regret capture
different kinds of deviations.

The corresponding deviation classes are summarized in
Table~\ref{tab:regret-deviations}.

\begin{table}[ht]
    \centering
    \caption{Deviation classes associated with the regret notions.}
    \label{tab:regret-deviations}
    \begin{tabular}{ll}
        \toprule
        Regret notion & Hindsight deviation \\
        \midrule
        External
        & Replace all realized actions by one fixed action \\
        Internal
        & Replace one action \(i\) by another action \(j\) \\
        Swap
        & Choose a replacement separately for every action \\
        \bottomrule
    \end{tabular}
\end{table}

For a single realized action sequence, the experiments compute cumulative
action regret \(R_T\) and average action regret \(R_T/T\). A learner is
no-regret with respect to a given deviation class if
\[
    \limsup_{T\rightarrow\infty}\frac{R_T}{T}\leq 0.
\]
For internal and swap regret as defined above, the identity replacement makes
the regret nonnegative, so this reduces to \(R_T/T\rightarrow0\).

For comparison, replacing
\(\mathbf{1}\{I_t=i\}\) in \(G_T(i,j)\) by \(p_t(i)\) gives
\[
    \overline{G}_T(i,j)
    =
    \sum_{t=1}^{T}
    p_t(i)\bigl(r_t(j)-r_t(i)\bigr).
\]
This is the corresponding \emph{distribution-regret} construction in the
terminology of \citet{greenwald_bounds_2006}: it evaluates the mixed strategy
rather than the sampled action sequence. Theoretical guarantees therefore need
not refer to the same regret object or probabilistic criterion: the literature
distinguishes action from distribution regret \citep{greenwald_bounds_2006}
and pseudo-regret from stronger high-probability guarantees
\citep{huang_near-optimal_2023}. These distinctions are stated explicitly
where the algorithms are introduced in Chapter~\ref{chap:algorithms}.

The experiments in this thesis use action regret under both full-information
and bandit feedback, and report its replicate mean as an empirical estimate of
expected action regret. In the bandit experiments, the learner still receives
only \(r_t(I_t)\); rewards of actions that were not selected are retained by
the experimental environment solely for offline regret evaluation.


\section{Equilibrium Concepts}
\label{sec:equilibrium-concepts}

The repeated-game experiments require several standard equilibrium concepts.
A finite normal-form game consists of a set of players
\(N=\{1,\ldots,n\}\), finite action sets \(A_i\), and payoff functions
\(u_i:A_1\times\cdots\times A_n\rightarrow\mathbb{R}\).

Let \(a_t=(a_{1,t},\ldots,a_{n,t})\) be the realized joint action in round
\(t\). The empirical distribution of joint play after \(T\) rounds is
\[
    \widehat{\mu}_T(a)
    =
    \frac{1}{T}
    \sum_{t=1}^{T}
    \mathbf{1}\{a_t=a\}.
\]
This distribution is used directly in the equilibrium analysis of the
experiments.

\begin{definition}[Coarse correlated equilibrium]
A distribution \(\mu\) over joint actions is a
\emph{coarse correlated equilibrium} (CCE) if, for every player \(i\) and
every fixed alternative action \(a_i'\in A_i\),
\[
    \sum_a
    \mu(a)
    \bigl(
        u_i(a_i',a_{-i})-u_i(a)
    \bigr)
    \leq0.
\]
\end{definition}

A CCE therefore rules out profitable deviations to a single fixed action
chosen independently of the action drawn from \(\mu\). Replacing the right-hand
side \(0\) by \(\varepsilon\geq0\) gives the corresponding
\(\varepsilon\)-CCE notion \citep{leme_convergence_2024}.

\begin{definition}[Correlated equilibrium]
A distribution \(\mu\) over joint actions is a
\emph{correlated equilibrium} (CE) if, for every player \(i\) and every pair
of actions \(a_i,a_i'\in A_i\),
\[
    \sum_{a_{-i}}
    \mu(a_i,a_{-i})
    \bigl(
        u_i(a_i',a_{-i})
        -
        u_i(a_i,a_{-i})
    \bigr)
    \leq0.
\]
\end{definition}

Unlike CCE, a correlated equilibrium rules out deviations that depend on the
action assigned to the player. Replacing the right-hand side \(0\) by
\(\varepsilon\geq0\) defines an \(\varepsilon\)-CE. This parallels the
action-dependent replacement considered by internal and swap regret.
Correlated equilibrium was introduced by \citet{aumann_subjectivity_1974}.

Every correlated equilibrium is also a coarse correlated equilibrium, so
\(\mathrm{CE}(G)\subseteq\mathrm{CCE}(G)\)
\citep{leme_convergence_2024}.

The experiments later measure the distance of \(\widehat{\mu}_T\) to these
equilibrium sets; the precise numerical procedure is described in
Chapter~\ref{chap:methodology}.


\section{Regret and Equilibrium}
\label{sec:theoretical-connections}

The standard regret--equilibrium correspondence motivates the repeated-game
experiments. If every player has vanishing external regret, the empirical
distribution of joint play approaches the coarse correlated-equilibrium set
\citep{chen_hedging_2020}. If every player has vanishing internal or swap
regret, the corresponding empirical distribution approaches the
correlated-equilibrium set \citep{hart_simple_2000,blum_external_2007}. These
standard implications are used later only as qualitative reference points for
the empirical CE and CCE measurements.

A stronger result is known in a more specialized self-play setting.
\citet{leme_convergence_2024} show that, for almost all symmetric two-player
zero-sum games, if both players use the same no-swap-regret dynamics with
symmetric initialization, then all but a vanishing fraction of the
mixed-strategy iterates are close to a symmetric Nash equilibrium. The
experiments in this thesis do not test this frequent-iterate result directly;
they measure regret and the time-averaged joint distribution's distance to CE
and CCE.

The next chapter introduces the concrete learning algorithms implemented in
the empirical framework and summarizes their theoretical regret guarantees.
